% =============================================================
% ERPM'27 (Empirical Research in Process Mining @ ICPM'27) — 12pp LNBIP
% Cut from paper/paper-v2.tex at commit c489460 (the reviewed long version).
% Bibliography: ../references.bib
%
% RULE OBSERVED IN THIS FILE: every number appears either in paper-v2.tex
% at c489460 or in a committed artefact under strata/outputs/. The artefact
% behind each is named in a comment beside it.
%
% Sections 6 and 7 are PLACEHOLDERS by instruction and are not drafted here.
% =============================================================

\documentclass[runningheads]{llncs}

\usepackage[T1]{fontenc}
\usepackage{microtype}
\usepackage{booktabs}
\usepackage{url}

\begin{document}

\title{Process Mining Did Not Detect the Faults}
\subtitle{Four measured results from an application to HVAC, and the protocol
that refuted the hypothesis behind it}

%TODO co-authors: supervisors to be added as \inst{1} co-authors pending
%     confirmation of authorship order (A. Aighobahi, M. Sharma).
\author{Yassh Singh}
\authorrunning{Y. Singh}
\institute{Thompson Rivers University, Kamloops, BC, Canada
%TODO contact e-mail
}

\maketitle

\begin{abstract}
We report a negative result and the protocol that produced it. Discovering
Petri nets from fault-free operation of three building HVAC systems, we tested
whether conformance against them detects seeded faults. Over all 73 seeded
scenarios no fault is detected only by the two conformance channels; removing
them changes no scorecard and lowers one system's false-alarm rate from four
days in 96 to one. Three further measurements support that reading: a measured
circularity bound, 1.78\% against 5.24\% recall at one threshold, obtained by
partitioning the event alphabet; a reversal probe showing that two of the three
discovered nets carry no ordering information in their fitness; and
alignment-based support that equalled the raw log count exactly at the unit
stratum, the net's reading adding nothing over counting. We report the
protocol --- pre-registration with binding falsifiers, adversarial
recomputation, tests that pin disclosures --- as what made the null checkable.

\keywords{Process mining \and Conformance checking \and Event abstraction \and
Negative results \and Reproducibility \and Cyber-physical systems.}
\end{abstract}

% =============================================================
\section{Introduction: an empty quadrant, and what we found in it}
\label{sec:intro}

Process mining has spread steadily from business processes into
cyber-physical and sensor-driven domains, and the spread has been reported
as a success: discovery plus conformance checking works on robotic arms, on
industrial IoT streams, on pasteurisation lines governed by an explicit
regulatory standard. One domain of that kind is missing from the literature.
Building heating, ventilation and air-conditioning (HVAC) operation is a
continuous-flow, standard-governed, heavily instrumented process, and it does
not appear. The nearest systematic review, of 36 process-mining studies on
sensor data, classifies fourteen industrial applications, eleven smart-home
applications, seven healthcare, two smart-office, two smart-product and one
other, and records no building HVAC~\cite{brzychczy2025pmsensorreview}. The
absence is corroborated from the other side: a survey of 211 AI-based HVAC
fault-detection studies published between 2013 and 2023 has no process-mining
entry at any leaf of its method taxonomy~\cite{bi2024aihvacreview}, a review of
47 process-mining studies in Industry 4.0 covers manufacturing, logistics,
mining and alarm management but no building
applications~\cite{akhramovich2024pmindustry40}, where process mining has
reached the built environment it is for construction and facility-management
workflows with cases as projects~\cite{martinezlagunas2024constructionpm}, and
the canonical fault-detection taxonomy for the domain does not list the method
at all~\cite{granderson2017fddtoolsurvey}.

The quadrant looked worth filling for a specific reason. The largest empirical
characterisation of HVAC faults to date, across more than 60{,}000 pieces of
equipment in 317 commercial buildings, finds that of the ten most common air
handler faults, seven are behavioural rather than
component-level~\cite{crowe2023faultprevalence}: schedules that run at the
wrong hours, manual overrides never released, setpoint resets that are skipped,
economiser sequences that execute their steps out of order. These are
deviations in the ordering of operating states. Threshold alarms and
hand-written rule libraries, which is what buildings actually run, fire on a
variable leaving a band or on a named pattern being matched, and neither
mechanism sees a sequence executed in the wrong order unless someone wrote a
rule for that sequence in advance. The structural argument is
straightforward, and it was our hypothesis: a model of the order in which a
building moves between operating states should have an advantage on exactly the
fault classes that dominate the field. We pre-registered it, with the
observation that would refute it, and built the pipeline to test it.

It was refuted. This paper reports the refutation and the three measurements
that explain it, because a workshop on empirical research in process mining is
the right venue for a result of this shape and because the field has almost no
published record of what happens when the method is carried into a new domain
and does not work. A precedent exists in another cyber-physical domain: Zhong
and Lisitsa report, in a preprint, that na\"{\i}ve process mining on network
event logs was ineffective for intrusion detection~\cite{zhong2022pmids}. Ours
is the second data point we are aware of, and unlike the first it is
accompanied by artefacts from which each number regenerates.

Four results follow, in the order in which the protocol produced them. First, a
\emph{measured circularity bound}: partitioning the event alphabet so that
events already encoding engineering suspicion are withheld from discovery
lowers the discovered model's apparent recall from 5.24\% to 1.78\% at the same
threshold on the same split, a 3.5-percentage-point bound on model
self-confirmation on a sensor-derived log (Section~\ref{sec:abstraction}).
Second, the \emph{reversal probe}, a one-line test of whether a discovered net
encodes order at all: fitness falls from 0.8861 to 0.5526 on one net and does
not move on the other two (Section~\ref{sec:discovery}). Third, an
\emph{artefact-verified null}: over 73 seeded scenarios, no fault is detected
only by the two conformance channels, and removing those channels changes no
scorecard (Section~\ref{sec:null}). Fourth, at the unit stratum the
\emph{alignment-based support equalled the raw log count exactly}, so the net's
reading added nothing over counting, while the device stratum ordered the two
comparable buildings the opposite way (Section~\ref{sec:discovery}). The HVAC
domain is the setting for these results; the results themselves are about
process mining.

% =============================================================
\section{Related work}
\label{sec:related}

\paragraph{Conformance as an anomaly detector.}
Using conformance checking directly as a trace-level anomaly detector is an
established line: Bezerra and Wainer score traces against a discovered model to
flag anomalous ones~\cite{bezerra2013anomalytraces}, on conformance
foundations laid by Rozinat and van der
Aalst~\cite{rozinat2008conformance,vanderaalst2012replaying,carmona2018conformancechecking}.
The line then moved elsewhere. Nolle and colleagues' BINet and its successors
are neural sequence models trained on event logs, and they are reported to
outperform conformance-based baselines on standard log
corpora~\cite{nolle2022binet}. Our null belongs with that migration rather than
against it: the conformance channels described here are the nearest descendant
of the earlier approach, and finding that they contribute no detection a
calibrated statistical channel does not already contribute is consistent with
why the field moved.

\paragraph{Event abstraction.}
Turning continuous sensor streams into discrete activity sequences is the
binding technical step for any sensor-domain application. We follow the
methodology of Janssen and colleagues for deriving event logs from sensor
data~\cite{janssen2021sensoreventlog} within the abstraction taxonomy of van
Zelst and colleagues~\cite{vanzelst2021eventabstraction}.
Section~\ref{sec:abstraction} reports the one decision in that step that
carried a measurable consequence.

\paragraph{Leakage.}
That decision is a leakage decision. Weytjens and De Weerdt catalogue data
leakage in process-mining experiments and show how easily a pipeline can
evaluate itself~\cite{weytjens2022dataleakage}. The alphabet partition of
Section~\ref{sec:abstraction} is a leakage control, and its cost --- what the
model loses when it is no longer allowed to recognise the analyst's own rules
--- is the first of our four results.

\paragraph{Precision, and why we did not use it.}
Whether a discovered net constrains order or merely permits a set of activities
is the question precision metrics are meant to
answer~\cite{munozgama2010precision,buijs2012quality}. Precision is, however, a
contested measurement: Tax and colleagues catalogue its
imprecisions~\cite{tax2018imprecisions}, and vanden Broucke and colleagues show
how ill-posed it is on logs without negative
examples~\cite{vandenbroucke2014negativeevents}, which sensor logs never have.
We therefore replaced the metric with a direct probe
(Section~\ref{sec:discovery}).

\paragraph{Hypothesis testing in process mining.}
Pitsch and colleagues argue that process-mining research should adopt explicit
hypothesis testing~\cite{pitsch2025hypothesistesting}. The protocol behind this
paper is one attempt to practise that, pre-registration and binding falsifiers
included; Section~\ref{sec:protocol} describes the part of it that is
transferable.

\paragraph{HVAC fault detection, in three sentences.}
The dominant research method in this domain is supervised machine learning,
surveyed across 211 studies whose reported accuracies have saturated at
97--99\% on the standard benchmarks~\cite{bi2024aihvacreview}, with the recent
shift toward interpretability implemented as post-hoc feature attribution on a
black-box model. What buildings actually deploy is rule-based: commercial
fault-detection products match libraries of pre-encoded fault
patterns~\cite{granderson2017fddtoolsurvey}, and the strongest research
programme in that style field-tested seven hand-coded correction algorithms
across 225 equipment items~\cite{pritoni2022faultcorrection}. An audit of 65
primary machine-learning studies in the domain found two that shared code
links, one of them broken, and 72\% that did not state whether their data was
public~\cite{mukhtar2025reproducibility}; evaluation here uses the public,
DOI-citable LBNL benchmark family~\cite{granderson2020lbnlfdddata}.

\paragraph{Hand-designed Petri nets are a different class.}
Petri nets have modelled HVAC and building systems for two decades --- hybrid
system modelling~\cite{villani2004hybridpetri}, building-automation
monitoring~\cite{gomes2007petribas}, chiller fault
diagnosis~\cite{xu2015chillerpetri} --- but in every case an engineer authored
the net, there is no event log and no discovery algorithm, and that body of
work should not be read as a precedent for discovering a net from a log of
building operation.

% =============================================================
\section{Event abstraction, the alphabet wall, and a measured circularity
bound}
\label{sec:abstraction}

\subsection{The abstraction, and the case notion problem}

HVAC operation offers no natural case boundary. A business process has
instances; a building runs continuously, with concurrent equipment at several
granularities and signals that are sampled rather than emitted. We take the
calendar day as the case, which is the shortest unit over which the control
sequences of interest --- start-up, occupancy, economiser transitions,
night cycling, shutdown --- complete. Continuous streams are converted to
events by control-sequence rules grounded in ASHRAE Guideline 36, following
Janssen and colleagues~\cite{janssen2021sensoreventlog}. The canonical
alphabet has ten activities in five complementary pairs: system started and
stopped, night cycle started and ended, heating active and inactive, cooling
active and inactive, economiser window entered and
exited.\footnote{\texttt{outputs/grammar\_results.json},
\texttt{canonical\_alphabet}.} Not every system emits all ten: the single-duct
air handler's configuration yields six activities, and both fan-powered-unit
configurations yield ten.

\subsection{The alphabet wall}

The events are then partitioned into two alphabets that are never mixed.
\emph{State events} are neutral facts about what the equipment did --- the fan
started, cooling became active --- and are the only events admitted to model
discovery. \emph{Fault-signature events} already encode engineering suspicion:
they exist because an analyst decided that a particular configuration of
measurements is worth naming. They are withheld from discovery entirely.

The partition exists because of circularity, which is the process-mining form
of the leakage problem catalogued by Weytjens and
De~Weerdt~\cite{weytjens2022dataleakage}. A model discovered over an alphabet
that already names faults will later ``detect'' those faults, but what it has
recognised is the analyst's own rule, re-encoded as a transition and then read
back out through alignment cost. The detection is real in the sense that the
number is correctly computed and vacuous in the sense that it measures the
pipeline against itself. The literature warns about this in the abstract; we
are not aware of a prior quantification of it on a sensor-derived log.

\subsection{The bound}

Quantifying it needs only one comparison. We removed the partition and re-ran
the discovery and scoring pipeline on the single-duct air-handler system,
holding everything else fixed: the same alignment threshold, the same 248-day
training and 87-day holdout split, the same 4{,}540 fault days scored. The
model channel's apparent recall rises from 1.78\% to 5.24\% --- 81 detected
fault days against 238 --- purely because the discovered net is now allowed to
contain the activities that name the faults.\footnote{Partitioned run:
\texttt{outputs/benchmark\_v3.json}. Unpartitioned predecessor:
\texttt{outputs/benchmark\_v2\_processheal\_v1.json}. Both report threshold
$0.333$, \texttt{train\_days} 248, \texttt{holdout\_days} 87, and
\texttt{pos\_days} 4{,}540; the \texttt{structure} channel's true positives are
81 and 238 respectively.} The difference, 3.5 percentage points, is a measured
bound on model self-confirmation for this system and this alphabet: it is how
much of the model channel's performance was the model recognising work that had
already been done for it.

Three things should be said about what this number is. It is a bound on one
system, and the size of the effect must depend on how much of the alphabet
encodes suspicion; a pipeline with a larger signature alphabet would leak more.
It is a lower bound on the harm in one direction only --- it measures inflated
recall, not the more insidious possibility that a signature-contaminated net
becomes \emph{better} at the faults the analyst already understood and no
better at anything else. And it is the smaller half of a ratio whose larger
half is that even the inflated 5.24\% is not a useful detection rate; the
partition removed two thirds of a capability that was not, in absolute terms,
worth having. That the pipeline's most defensible methodological decision
raised no result and lowered one is the honest summary.

\subsection{Pinning the partition}

A partition enforced by convention decays. The separation is therefore held by
a tested invariant: injecting fault-signature events into an event log must not
change any day's fitness against a discovered net. The test fails if a future
edit lets signature events reach the discovery path, and it fails at build time
rather than in a reviewer's reading. We recommend the pattern to anyone
building a process-mining pipeline over engineered events, independently of
whether the rest of this paper's results replicate: the alphabet is where the
leakage lives, and a prose statement that two alphabets are separate is not a
control.

% =============================================================
\section{Discovery at three strata, and the reversal probe}
\label{sec:discovery}

\subsection{Three strata}

Petri nets are discovered from the state-event log with the inductive
miner~\cite{leemans2013inductiveminer,leemans2014inductiveminerinfrequent} as
implemented in pm4py~\cite{berti2023pm4py}, at three granularities. The
\emph{device} stratum models an individual actuator's lifecycle, with one case
per device per day. The \emph{unit} stratum models the air handler's daily
operating cycle. A \emph{cross-system grammar} maps each system's alphabet onto
the canonical ten activities so that a net discovered on one building can be
replayed against another building's log. All discovery uses fault-free data
only.

\subsection{The reversal probe}

Whether such a net encodes temporal order, or merely licenses a set of
activities in any arrangement, is the question precision metrics
address~\cite{munozgama2010precision,buijs2012quality}, and it is the question
they answer least reliably on logs with no negative
examples~\cite{tax2018imprecisions,vandenbroucke2014negativeevents}. We
therefore use a direct probe. Reverse every trace in the log, re-score it
against the same net, and report the drop in mean fitness. A net that
constrains order cannot fit a reversed log as well as the original; a net whose
fitness is a function of the activity multiset alone will not notice.

The three nets separate cleanly.\footnote{\texttt{outputs/grammar\_results.json},
\texttt{reversal\_probe}.} The single-duct net falls from 0.8861 to 0.5526,
a drop of 0.3335: it encodes order. The series fan-powered
net scores 0.9533 forwards and 0.9533 backwards, identical to four decimal
places. The parallel fan-powered net moves from 0.9021 to 0.9026 --- five
ten-thousandths, in the wrong direction. Two of the three nets carry no
ordering information in their fitness at all, and their high absolute fitness
is a statement about which activities occurred, not about when.

This matters beyond the diagnosis. A fitness of 0.95 reads as an excellent
model; it is the number one would quote. The probe costs one pass over the log
and one re-scoring, needs no negative examples and no reference model, and it
is what told us that two thirds of our model estimate was not doing the thing
the method is chosen for. We recommend it as a routine reporting item
alongside fitness in any sensor-domain application, and we would have caught
the problem two phases earlier had we run it first.

\subsection{Shared structure exists only where order does}

The cross-system grammar gives six off-diagonal cells, each a net discovered on
one system replayed against another system's log, with a null built by
replaying five shuffled versions of the same log through the same net and
taking the mean. Two cells clear their null: fan-powered-unit logs outscore
their own shuffled controls on the single-duct net by 0.18 and 0.15 in mean
fitness. The remaining four are null, with order signals between $-0.005$ and
$0.000$.\footnote{\texttt{outputs/grammar\_results.json}, \texttt{M1\_null}:
\texttt{order\_signal} $0.1828$ and $0.1537$ for the two positive cells;
$-0.0047$, $-0.0007$, $-0.0040$ and $-0.0007$ for the four null ones.} The
pattern is not arbitrary: both positive cells are cells in which the
\emph{single-duct} net is the model, and the single-duct net is the one the
reversal probe found order-sensitive. Real shared behavioural structure
between two buildings is detectable only through a net that encodes structure
in the first place, which turns the reversal probe from a diagnostic into a
precondition for cross-system claims.

\subsection{The net's reading added nothing over counting}

A fourth measurement belongs here although it was computed for the test
reported in Section~\ref{sec:transfer}. For each system we asked how often a
named activity is genuinely supported by the model, counting the days on which
that activity is replayed as a synchronous move against the discovered unit
net. On every system the answer equalled the number of days on which the
activity simply appears in the raw log: 0 against 0 on the single-duct system,
166 against 166 on the parallel unit and 357 against 357 on the
series unit.\footnote{\texttt{outputs/discovery\_predicts\_transfer.json},
\texttt{Q}: \texttt{sync\_days} against
\texttt{\_aux.occupied\_days\_whose\_log\_contains\_heating}.} At the unit
stratum the alignment added nothing that counting occurrences would not have
given, exactly. The device stratum, whose net is pooled across the four
terminal units of each building, disagrees with the unit stratum about which of
the two comparable buildings has the better-supported behaviour: supports of
$0.455$ and $0.978$ at the unit stratum against $0.437$ and $0.014$ at the
device stratum, an inversion. A quantity that is reproducible by counting on
one stratum and reverses on another is not yet a measurement, and
Section~\ref{sec:transfer} reports what happened when we tried to build a test
on it.

% =============================================================
\section{The conformance null over 73 seeded scenarios}
\label{sec:null}

\subsection{Eight channels, two of them process mining}

Evaluation used three public LBNL fault-detection
datasets~\cite{granderson2020lbnlfdddata}: a single-duct air-handling unit and
parallel and series fan-powered-unit configurations, 73 seeded fault scenarios
in total, scored at day level with the scenario as the unit of analysis
following the established evaluation
protocols~\cite{yuill2013fddprotocol,frank2018fddevaluation}. Eight channels
score each day independently: (1) \emph{rules}, physics rules naming direct
contradictions between measurements; (2) \emph{residual}, temperature-residual
bands; (3) \emph{model conformance} and (4) \emph{device conformance},
alignment cost against the discovered unit and device nets respectively ---
these two, and only these two, are the process-mining channels; (5)
\emph{absence}, a scheduled activity that did not occur; (6) \emph{frequency}
and (7) \emph{oscillation}, statistical bands on activity counts and reversals;
and (8) a seasonal \emph{rate} channel that corroborates but may not alarm
alone. Every channel is calibrated on the fault-free year only and may report
only what exceeds its own measured noise floor, tested as a binomial against
that channel's own holdout rate at $p < 10^{-3}$. This design is what makes the
ablation meaningful: the conformance channels are not competing against an
uncalibrated strawman, and each of the other channels had to clear the same
gate.

\subsection{The ablation}

Because the committed scorecards record, per scenario, which significance-gated
channels were credited with the detection, the ablation is a read rather than a
re-run. The result is in Table~\ref{tab:ablation}. No scenario on any system is
detected only by the two conformance channels. They are credited with none of
the 14 single-duct detections and none of the 23 parallel-unit detections, and
they appear in 7 of the 24 series-unit detections but never alone. Removing
both channels from the detector therefore leaves every scorecard unchanged.

\begin{table}[t]
\centering
\caption{The conformance ablation, read from the committed scorecards
(\texttt{outputs/benchmark\_v6\_\{sdahu,pfpu,sfpu\}.json},
\texttt{meaningful\_channels}). ``Credited'' counts detections in which a
conformance channel was among the gated channels; ``sole detector'' counts
detections in which conformance was the only one. The single-duct count is the
naive 14 of 14; one of those detections is later adjudicated out on
dataset-provenance grounds, which changes no cell of this table because
conformance is credited with none of them.}
\label{tab:ablation}
\begin{tabular}{lrrrr}
\toprule
System & Scenarios & Detected & Conformance credited & Conformance sole detector \\
\midrule
Single-duct AHU      & 14 & 14 & 0 & 0 \\
Parallel fan-powered & 30 & 23 & 0 & 0 \\
Series fan-powered   & 29 & 24 & 7 & 0 \\
\midrule
Total                & 73 & 61 & 7 & 0 \\
\bottomrule
\end{tabular}
\end{table}

Removing them is not merely free. On the series unit the two conformance
channels account for all four false-alarm days on the 96-day fault-free
holdout; deleting them eliminates three of the four outright, the fourth being
independently flagged by the frequency channel, so that system's false-alarm
rate falls from 4 in 96 to 1 in 96 at no detection
cost.\footnote{\texttt{outputs/union\_fpr\_sfpu.json}: model channel
\texttt{holdout\_fp\_dates} \{2018-11-29\}, device channel \{2018-01-25,
2018-01-30, 2018-12-28\}, frequency channel \{2018-11-29\};
\texttt{union\_all8} 4 days.} The process-mining channels' entire measurable
contribution to this detector was three false alarms.

\subsection{The matched-rule control, and what it says the models were doing}

A second control was run to establish what discovery contributes as a
\emph{design-time} activity, separately from what it contributes at run time.
For each channel we asked a domain expert to hand-write the rule the channel's
behaviour implied, and scored the rule on the same scenarios. Two of these
reproduce their channel exactly: 141 detections against 141 on the series unit
and 206 against 206 on the parallel unit, in both cases the frequency channel
on the unstable-room-temperature
scenario.\footnote{\texttt{outputs/matched\_rules\_\{sfpu,pfpu\}.json}: for
\texttt{SFPU\_RMTEMPUnstable}, \texttt{mr2} $=$ \texttt{freq} $= 141$; for
\texttt{PFPU\_RMTEMPUnstable}, $=206$.} This is the useful half of the finding.
Discovery located an invariant nobody had written down, and once it was
written down the automated channel was exactly redundant with it.

The original hypothesis --- that discovered models out-detect hand-written
rules --- had a falsifier attached, and the falsifier fired. On the series unit,
a one-line matched rule reading \emph{occupied workday and air-handler heating
never active} covered 192 fault days against the model channel's 135 on the
same scenario, with zero false alarms on 365 fault-free
days.\footnote{\texttt{outputs/matched\_rules\_sfpu.json},
\texttt{SFPU\_VAVDMPRStuck\_80\%}: \texttt{mr1} $=192$, \texttt{model} $=135$;
\texttt{healthy\_fp.mr1} $=0$.} The 135 should not be read without its
robustness counterpart: replacing the interpolated threshold with the most
conservative one leaves 22 to 26 days per scenario, so roughly four fifths of
the threshold count rests on a threshold interpolated from a single holdout
day.

The control also carries a retraction we owe the reader, because the project
originally read this number the wrong way round. Transplanted to a building
whose configuration does not support it, the same one-line rule fires on 231
fault-free days on the single-duct system and 124 on the parallel unit, and on
none at all on the series unit.\footnote{\texttt{healthy\_fp.mr1} in
\texttt{outputs/matched\_rules\_\{sdahu,pfpu,sfpu\}.json}: 231, 124, 0.} We
first read the 231 as evidence about what the models do and do not predict. It
is not. The single-duct system has no heating rule and no hot-water valve, so a
rule reading \emph{heating never active} fires on every occupied day for want
of a sensor, and no model was involved in producing that number. The only
informative contrast the transplant offers is the parallel unit against the
series unit, 124 against 0, and whether the discovered models predict that
contrast is the subject of Section~\ref{sec:transfer}.

\subsection{Reading the null}

Three readings are available and we prefer the third. The conformance channels
might be poorly implemented: possible, but they use standard alignment-based
replay~\cite{vanderaalst2012replaying,carmona2018conformancechecking} through
pm4py on nets discovered by a standard miner, and the failure is not one of
tuning --- they detect, they simply detect nothing that is not already detected.
The domain might be wrong for the method: partly, and the reversal probe says
which part, since two of the three nets do not encode order and could not have
supplied the order-sensitivity the hypothesis rested on. The third reading is
the one the audits converged on. Every genuine model-channel detection we
inspected was a \emph{count} deviation --- extra heating episodes, collapsed fan
cycles, a vanished daily rhythm --- surfacing through alignment cost, and once a
frequency channel existed that measured count deviations directly and
two-sidedly, the conformance channels had nothing left that was uniquely theirs.
Alignment cost was acting as a noisy, one-sided, expensive frequency detector.
If that reading generalises, it predicts that conformance-as-anomaly-detection
will underperform wherever the discovered net is order-permissive, which is
wherever the alphabet is loop-heavy --- and sensor-derived alphabets, built from
complementary activity pairs that cycle many times a day, are loop-heavy almost
by construction.

% =============================================================
\section{A pre-registered test that never had power}
\label{sec:transfer}

% PLACEHOLDER: supplied by the lead after Amendment 3 is folded into the
% source manuscript.

% =============================================================
\section{Tests that pin refutations and disclosures}
\label{sec:protocol}

% PLACEHOLDER

% =============================================================
\section{Limitations, and what would make this replicable}
\label{sec:limitations}

\subsection{Limitations}

Every result here comes from simulation, and simulation flatters. Noise-free
years make zero-false-positive thresholds cheap, bands razor-thin and silence
easy. The direction of the bias is worth stating precisely for the null: a
noise-free training year should, if anything, \emph{help} a discovered net,
because the fault-free behaviour it must generalise from is unusually clean. We
do not think the null is an artefact of simulation, but we cannot exclude that
real, noisier logs would change the relative standing of conformance against
the statistical channels in either direction.

Three buildings is the harder limit, and it is a limit on the unit of
replication rather than on the number of days scored. There are 73 scenarios
and, on the single-duct system alone, 4{,}540 scored fault days, but a building
is the unit across which a discovered model must transfer, and there are three
of them, two of which are variants of one equipment family; all three share the
same simulated weather-year family, which is favourable to transfer. Any cross-building claim in this paper therefore rests on an $n$ small
enough that the pre-registered test built on it turned out to have no
attainable significance under any data (Section~\ref{sec:transfer}). A powered
version needs more buildings, more than two levels of the predictor, and an
outcome that is not fitted on the same year it is scored on.

The noise-floor gate applies no correction across channels and scenarios. On
the order of 584 binomial tests run at $\alpha = 10^{-3}$ --- eight channels by
73 scenarios, with the device channel additionally run per terminal unit ---
giving a family-wise error near 0.44 under independence and an expectation of
about 0.6 false cells across the study; per-day observations within a scenario
are also autocorrelated. Separately, the two conformance channels'
false-positive rows are calibration targets rather than out-of-sample
measurements, because their thresholds are quantile-fit on the same holdout
days those rows report. Both caveats cut against the conformance channels
looking \emph{better} than they are, which is the direction that matters for a
null.

Finally, the evaluation produced five machine-verified defects in the LBNL
benchmark datasets, including one in which the fault-free file was simulated on
a different configuration branch from every fault file; these are reported with
their reproducing evidence in a companion paper (in preparation), and one of
them is why Table~\ref{tab:ablation}'s single-duct row carries an adjudication
note.

\subsection{Disclosure: AI agents as part of the method}

AI coding agents were used for implementation, analysis scripting and
adversarial recomputation of results, and one step depended on them
methodologically rather than incidentally: the sealed quantity behind
Section~\ref{sec:transfer} was computed by an agent that was denied access to
the rule-firing artefacts it would be scored against, which is how the
computation was kept blind. All experimental designs, pre-registrations and
claims were authored and approved by the author.

\subsection{Conclusion}

We carried event-log process mining into a domain it had not been applied to,
with a pre-registered hypothesis and a falsifier attached, and the falsifier
fired. What survives as a working detector contains no process mining: the
channels that detect are physics rules, calibrated residual bands and
statistical bands on activity counts and reversals, and removing both
conformance channels changes no scorecard and removes three false alarms. The
honest statement of the outcome is that discovery earned its place at design
time, by locating an invariant that an expert could then write down as a rule,
and lost it at run time, exactly and measurably, to the rule it had suggested.

The three supporting measurements are what we would most like replicated,
because each is cheap and none is specific to buildings. The alphabet partition
carries a measured cost, 1.78\% against 5.24\% at one threshold on one split,
and we know of no prior quantification of circularity on a sensor-derived log;
the number should be measured again on other pipelines, and we expect it to be
larger where more of the alphabet encodes suspicion. The reversal probe costs
one re-scoring and told us that two nets with fitness above 0.90 carried no
order information whatever, which no fitness figure and no precision metric
communicated. And alignment-based support that reproduces a raw log count
exactly is a warning that the model is not being read, merely consulted. A
workshop on empirical research in process mining is the right place to say that
we ran the method properly, measured what it contributed, and found the
contribution to be zero; the contribution of the protocol that established this
is not zero.

\subsubsection*{Acknowledgements.}
The author thanks Dr.\ Anthony Aighobahi and Dr.\ Mridula Sharma for
supervision. AI coding assistants were used for implementation, analysis
scripting and adversarial auditing, and a sealed computation by an agent
blinded to the outcome artefacts is part of the method reported in
Section~\ref{sec:transfer}; all designs, pre-registrations and claims were
reviewed and approved by the author.

\bibliographystyle{splncs04}
\bibliography{../references}

\end{document}
